\documentclass[10pt,a4paper]{amsart}
\usepackage{amsmath,amssymb,amsthm,graphicx}
% Retained minimal preamble for the Stacks Project chapter varieties.tex.
% Verbatim excerpt of the chapter preamble: Xy-pic setup (native preamble lines 14-18), the native theorem
% declarations with the mathematical macros (native preamble lines 154-217), and the project own no-op
% definitions of the reference/slogan/history display environments (native wrapper.tex lines 1-3) that the
% counted statement uses for its slogan. Copied mechanically; nothing added, removed or redefined.
% For commutative diagrams we use Xy-pic
\usepackage[all]{xy}

% We use 2cell for 2-commutative diagrams.
\xyoption{2cell}
% Theorem environments.
%
\theoremstyle{plain}
\newtheorem{theorem}[subsection]{Theorem}
\newtheorem{proposition}[subsection]{Proposition}
\newtheorem{lemma}[subsection]{Lemma}

\theoremstyle{definition}
\newtheorem{definition}[subsection]{Definition}
\newtheorem{example}[subsection]{Example}
\newtheorem{exercise}[subsection]{Exercise}
\newtheorem{situation}[subsection]{Situation}

\theoremstyle{remark}
\newtheorem{remark}[subsection]{Remark}
\newtheorem{remarks}[subsection]{Remarks}

\numberwithin{equation}{subsection}

% Macros
%
\def\lim{\mathop{\mathrm{lim}}\nolimits}
\def\colim{\mathop{\mathrm{colim}}\nolimits}
\def\Spec{\mathop{\mathrm{Spec}}}
\def\Hom{\mathop{\mathrm{Hom}}\nolimits}
\def\Ext{\mathop{\mathrm{Ext}}\nolimits}
\def\SheafHom{\mathop{\mathcal{H}\!\mathit{om}}\nolimits}
\def\SheafExt{\mathop{\mathcal{E}\!\mathit{xt}}\nolimits}
\def\Sch{\mathit{Sch}}
\def\Mor{\mathop{\mathrm{Mor}}\nolimits}
\def\Ob{\mathop{\mathrm{Ob}}\nolimits}
\def\Sh{\mathop{\mathit{Sh}}\nolimits}
\def\NL{\mathop{N\!L}\nolimits}
\def\CH{\mathop{\mathrm{CH}}\nolimits}
\def\proetale{{pro\text{-}\acute{e}tale}}
\def\etale{{\acute{e}tale}}
\def\QCoh{\mathit{QCoh}}
\def\Ker{\mathop{\mathrm{Ker}}}
\def\Im{\mathop{\mathrm{Im}}}
\def\Coker{\mathop{\mathrm{Coker}}}
\def\Coim{\mathop{\mathrm{Coim}}}

% Boxtimes
%
\DeclareMathSymbol{\boxtimes}{\mathbin}{AMSa}{"02}

%
% Macros for moduli stacks/spaces
%
\def\QCohstack{\mathcal{QC}\!\mathit{oh}}
\def\Cohstack{\mathcal{C}\!\mathit{oh}}
\def\Spacesstack{\mathcal{S}\!\mathit{paces}}
\def\Quotfunctor{\mathrm{Quot}}
\def\Hilbfunctor{\mathrm{Hilb}}
\def\Curvesstack{\mathcal{C}\!\mathit{urves}}
\def\Polarizedstack{\mathcal{P}\!\mathit{olarized}}
\def\Complexesstack{\mathcal{C}\!\mathit{omplexes}}
% \Pic is the operator that assigns to X its picard group, usage \Pic(X)
% \Picardstack_{X/B} denotes the Picard stack of X over B
% \Picardfunctor_{X/B} denotes the Picard functor of X over B
\def\Pic{\mathop{\mathrm{Pic}}\nolimits}
\def\Picardstack{\mathcal{P}\!\mathit{ic}}
\def\Picardfunctor{\mathrm{Pic}}
\def\Deformationcategory{\mathcal{D}\!\mathit{ef}}
\newenvironment{reference}{}{}
\newenvironment{slogan}{}{}
\newenvironment{history}{}{}
\usepackage{cleveref}
\crefname{originalreference}{original reference}{original references}
\Crefname{originalreference}{Original reference}{Original references}
\makeatletter
\newlabel{situation-glue}{{09MY}{source}{}{sourcepointer-1}{}}
\newlabel{situation-glue@cref}{{[originalreference][0][]09MY}{[1][0][]source}{}{}{}}
\makeatother
\title{Glueing a Noetherian semi-local domain of dimension one with a discrete valuation ring inside a common field}
\author{The Stacks Project Authors}
\date{stacks-project:varieties.tex:local-snapshot-20261011; GNU FDL 1.2 or later}
\begin{document}
\maketitle
\noindent\textit{Editorial note.} All delivered mathematical text is the author's own, copied verbatim from the pinned Stacks Project chapter file varieties.tex; no statement or proof step was rewritten, reordered or removed.\par\medskip
\noindent\textit{Editorial note.} The counted span cites one cross-chapter situation label, recorded in assembly.external\_references with its canonical tag as a navigation pointer.\par\medskip
\noindent\textit{Editorial note.} The proof cites no bibliography item inside the counted span.\par\medskip
\noindent\textit{Allowed standard prerequisites (editorial).}\begin{itemize}
\item The glueing situation for semi-local rings. Stacks Project, Varieties, Situation situation-glue (tag 09MY). Supplies the ambient hypotheses on A, B, K and R that the two case analyses use.
\end{itemize}\medskip
\ifdefined\hypertarget\hypertarget{sourcepointer-1}{}\fi
\noindent\textit{Source reference.} Original reference 09MY: sources/stacks-chow/tags.txt, source lines 12196--12196. Cross-chapter standard prerequisite cited by the author inside the counted span; it is used as a pointer and its own proof is not part of this item.\par\medskip
\begin{lemma}
\label{lemma-semi-local}
In Situation \ref{situation-glue} assume
\begin{enumerate}
\item $A$ is a Noetherian semi-local domain of dimension $1$,
\item $B$ is a discrete valuation ring,
\end{enumerate}
Then we have the following two possibilities
\begin{enumerate}
\item[(a)] If $A^*$ is not contained in $R$, then
$\Spec(A) \to \Spec(R)$ and $\Spec(B) \to \Spec(R)$
are open immersions covering $\Spec(R)$ and $K = A \otimes_R B$.
\item[(b)] If $A^*$ is contained in $R$, then $B$ dominates one of
the local rings of $A$ at a maximal ideal and $A \otimes B \to K$
is not surjective.
\end{enumerate}
\end{lemma}

\begin{proof}
Assumption (a) implies there is a unit $u$ of $A$ whose image in $K$ lies in
the maximal ideal of $B$. Then $u$ is a nonzerodivisor of $R$ and for
every $a \in A$ there exists an $n$ such that $u^n a \in R$. It follows
that $A = R_u$.

\medskip\noindent
Let $\mathfrak m_A$ be the Jacobson radical of $A$. Let $x \in \mathfrak m_A$
be a nonzero element. Since $\dim(A) = 1$ we see that $K = A_x$.
After replacing $x$ by $x^n u^m$ for some $n \geq 1$ and
$m \in \mathbf{Z}$ we may assume $x$ maps to a unit of $B$.
We see that for every $b \in B$ we have that $x^nb$
in the image of $R$ for some $n$. Thus $B = R_x$.

\medskip\noindent
Let $z \in R$. If $z \not \in \mathfrak m_A$ and $z$ does not map to
an element of $\mathfrak m_B$, then $z$ is invertible.
Thus $x + u$ is invertible in $R$. Hence $\Spec(R) = D(x) \cup D(u)$.
We have seen above that $D(u) = \Spec(A)$ and $D(x) = \Spec(B)$.

\medskip\noindent
Case (b). If $x \in \mathfrak m_A$, then $1 + x$ is a unit and
hence $1 + x \in R$, i.e, $x \in R$. Thus we see that
$\mathfrak m_A \subset R \subset A$. In fact, in this case
$A$ is integral over $R$. Namely, write
$A/\mathfrak m_A = \kappa_1 \times \ldots \times \kappa_n$
as a product of fields. Say $x = (c_1, \ldots, c_r, 0, \ldots, 0)$
is an element with $c_i \not = 0$. Then
$$
x^2 - x(c_1, \ldots, c_r, 1, \ldots, 1)  = 0
$$
Since $R$ contains all units we see that $A/\mathfrak m_A$ is
integral over the image of $R$ in it, and hence $A$ is integral over $R$.
It follows that $R \subset A \subset B$ as $B$ is integrally closed.
Moreover, if $x \in \mathfrak m_A$ is nonzero, then
$K = A_x = \bigcup x^{-n}A = \bigcup x^{-n}R$.
Hence $x^{-1} \not \in B$, i.e., $x \in \mathfrak m_B$.
We conclude $\mathfrak m_A \subset \mathfrak m_B$.
Thus $A \cap \mathfrak m_B$ is a maximal ideal of $A$ thereby
finishing the proof.
\end{proof}

\end{document}
